\documentclass[11pt,notheorems,aspectratio=169,t]{beamer}

% Shared package and font setup for Tao Slides Beamer.
% Build with XeLaTeX so the Beamer template can mirror the PowerPoint typography.

\usepackage{fontspec}
\usefonttheme{professionalfonts}
\IfFileExists{/mnt/c/Windows/Fonts/msyh.ttc}{
  \setsansfont[
    Path=/mnt/c/Windows/Fonts/,
    BoldFont=msyhbd.ttc
  ]{msyh.ttc}
  \setmainfont[
    Path=/mnt/c/Windows/Fonts/,
    BoldFont=msyhbd.ttc
  ]{msyh.ttc}
}{
  \IfFontExistsTF{Microsoft YaHei}{
    \setsansfont{Microsoft YaHei}
    \setmainfont{Microsoft YaHei}
  }{
  \setsansfont{TeX Gyre Heros}
  \setmainfont{TeX Gyre Heros}
  }
}

\usepackage{styles/beamerthemetao}

\newcommand{\makepart}[1]{
  \part{#1}
  {
    \setbeamertemplate{footline}{}
    \begin{frame}
      \partpage
    \end{frame}
  }
}

\RequirePackage{booktabs}
\RequirePackage{colortbl}
\RequirePackage{ragged2e}
\RequirePackage{tabularx}
\RequirePackage{array}
\RequirePackage{multirow}
\RequirePackage{caption}
\RequirePackage{subcaption}
\RequirePackage{wrapfig}
\RequirePackage{pgfplots}
\RequirePackage{graphicx}
\RequirePackage{adjustbox}
\RequirePackage{environ}
\RequirePackage{textcomp}
\RequirePackage{amsmath}
\RequirePackage{amsthm}
\RequirePackage{mathtools}
\newcolumntype{Y}{>{\centering\arraybackslash}X}
\pgfplotsset{compat=1.18}

\DeclareMathOperator*{\argmax}{arg\,max}
\DeclareMathOperator*{\argmin}{arg\,min}
\setbeamertemplate{theorems}[numbered]

\theoremstyle{definition}
\newtheorem{fact}{Fact}[section]
\newtheorem{examp}{Example}[section]

\theoremstyle{plain}
\newtheorem{definition}{Definition}[section]
\newtheorem{proposition}{Proposition}
\newtheorem{theorem}{Theorem}
\newtheorem{assumption}{Assumption}

\providecommand{\H}{\mathscr{H}}
\providecommand{\E}{\mathbb{E}}


\title[Title]{Title}
\author{Name}
\institute{Institution}
\date{20XX-XX-XX}

\begin{document}

{
\setbeamertemplate{footline}{}
\begin{frame}
  \titlepage
\end{frame}
}
\addtocounter{framenumber}{-1}

\begin{frame}{Page Title}{Subtitle}
  \begin{itemize}
    \item This template creates concise and consistent Beamer presentations.
    \item The visual style uses a white canvas, navy titles, a title rule, and bottom-right page numbers.
    \item Font sizes: page titles use 36 pt, subtitles use 24 pt, and body text uses 20 pt.
    \item Text uses Microsoft YaHei when available, and formulas use Computer Modern.
  \end{itemize}
\end{frame}

\begin{frame}{Page Title}{Subtitle}
  \begin{itemize}
    \item Item 1
    \item Item 2
    \item Item 3
  \end{itemize}
\end{frame}

\begin{frame}{Page Title}{Subtitle}
  \begin{itemize}
    \item Item 1
      \begin{itemize}
        \item Normal
        \item \emphasize{Emphasis}
        \item \alert{Important}
      \end{itemize}
  \end{itemize}
\end{frame}

\begin{frame}{Page Title}{Subtitle}
  \taobody{Content}
\end{frame}

\begin{frame}{Page Title}{Subtitle}
  \begin{taotwocolumns}
    \begin{column}{0.48\textwidth}
      \begin{taotable}
        \setlength{\tabcolsep}{0.12cm}
        \begin{tabularx}{\linewidth}{>{\raggedright\arraybackslash}Xccc}
          \taotablehead{Plan} & \taotablehead{A} & \taotablehead{B} & \taotablehead{Note} \\
          \taotableheadrule
          Baseline & 0.82 & 12.4 ms & Ref. \\
          \taotablerowrule
          Method 1 & 0.89 & 10.8 ms & \emphasize{Best} \\
          \taotablerowrule
          Method 2 & 0.87 & 9.6 ms & \taoaccent{Review} \\
          \taotablerowrule
        \end{tabularx}
      \end{taotable}
    \end{column}
    \begin{column}{0.48\textwidth}
      \begin{center}
        \begin{tikzpicture}[x=1cm,y=1cm]
          \draw[taoTableRule,line width=0.6pt] (0,0) -- (12.6,0);
          \draw[taoTableRule,line width=0.6pt] (0,1.2) -- (12.6,1.2);
          \draw[taoTableRule,line width=0.6pt] (0,2.4) -- (12.6,2.4);
          \draw[taoTableRule,line width=0.6pt] (0,3.6) -- (12.6,3.6);
          \draw[taoNavy,line width=1.5pt] (0,0) -- (0,4.0);
          \draw[taoNavy,line width=1.5pt] (0,0) -- (12.6,0);
          \draw[taoNavy,line width=2.2pt]
            (0.7,1.0) -- (3.2,1.8) -- (5.7,2.3) -- (8.2,2.9) -- (11.4,3.4);
          \draw[taoRed,line width=2.2pt]
            (0.7,2.9) -- (3.2,2.4) -- (5.7,2.0) -- (8.2,1.6) -- (11.4,1.1);
          \node[anchor=west,font=\fontsize{16}{20}\selectfont,text=taoNavy] at (0.2,4.35) {Metric Trend};
          \node[anchor=west,font=\fontsize{14}{18}\selectfont,text=taoNavy] at (7.2,4.35) {Metric A};
          \node[anchor=west,font=\fontsize{14}{18}\selectfont,text=taoRed] at (9.6,4.35) {Metric B};
        \end{tikzpicture}
      \end{center}
    \end{column}
  \end{taotwocolumns}
\end{frame}

\begin{frame}{Page Title}{Subtitle}
  \taobody{%
    Mathematical formulas use the Computer Modern font.

    \begin{align*}
      \frac{\partial n}{\partial t}
        &= \frac{1}{q}\nabla \cdot \mathbf{J}_n + G_n - R_n, \\
      \frac{\partial p}{\partial t}
        &= -\frac{1}{q}\nabla \cdot \mathbf{J}_p + G_p - R_p.
    \end{align*}
  }
\end{frame}

\end{document}
